\documentclass[10pt,a4paper]{amsart}
\usepackage[margin=19mm]{geometry}
\usepackage{amsmath,amssymb,mathtools}
\usepackage[scr=rsfs,cal=euler]{mathalfa}
\usepackage{xfrac}
\usepackage[unicode,hidelinks,bookmarks=false]{hyperref}
\newcommand{\ie}{i.e.\ }
\newcommand{\cf}{cf.\ }
\newcommand{\resp}{respectively\ }
\newcommand{\Subsection}{\subsection}
\providecommand{\nobreakdash}{\nobreak\hbox{-}}
\setlength{\emergencystretch}{2em}
\multlinegap0pt
\usepackage{bm}
\usepackage{bbm}
\usepackage{dsfont}
\usepackage{xspace}
\usepackage{cancel}
\usepackage{mathtools}
\usepackage{amsthm}
\makeatletter
\@ifundefined{thm}{\newtheorem{thm}{Thm}}{}
\@ifundefined{lem}{\newtheorem{lem}[thm]{Lemma}}{}
\@ifundefined{prop}{\newtheorem{prop}[thm]{Proposition}}{}
\makeatother
\newcommand{\RR}{\mathbb{R}}
\newcommand{\divg}{\mathrm{div}}
\newcommand{\faskip}{\hspace{8pt}}
\newcommand{\loc}{\mathrm{loc}}
\newcommand{\supp}{\mathrm{supp}}
\title[On the shape of the positivity region for a free boundary pr]{On the shape of the positivity region for a free boundary problem describing cell polarization}
\author{Sebasti\'an Flores Sep\'ulveda, Barbara Niethammer, Juan J. L. Vel\'azquez}
\date{Source: arXiv:2605.03553v1 (CC BY 4.0)}
\begin{document}
\maketitle

\noindent\emph{Source and licence.} On the shape of the positivity region for a free boundary problem describing cell polarization. Version arXiv:2605.03553v1 (CC BY 4.0), distributed under the Creative Commons Attribution 4.0 International licence, as recorded in the arXiv licence metadata for this exact version and shown on the abstract page https://arxiv.org/abs/2605.03553v1. Full rights notice: licenses/arxiv-2605.03553-CC-BY-4.0.txt.\par\smallskip

\noindent\textit{Editorial scope.} Counted unit: the Prop whose source label is \texttt{prop:exist\_noncoercive} in the article (source0.tex lines 1347--1350, source label \texttt{prop:exist\_noncoercive}), delivered together with the article's own complete original proof of it (source0.tex lines 1351--1404, one \texttt{proof} environment of 54 lines). No mathematics is written, repaired or re-derived: statement and proof are the article's own text line for line. Required context retained uncounted: lem:mp (lines 279--286); eq:obstacle\_f (lines 390--396); prop:existence (lines 399--402); prop:nondeg\_q\_bound (lines 510--513); eq:obst\_noncoercive (lines 1330--1336). Every definition, display and label of those lines is retained unchanged. Omitted from this package and not redistributed: 1--278, 287--389, 397--398, 403--509, 514--1329, 1337--1346, 1405--1672. Nothing inside a retained excerpt is omitted or altered: every retained body is delivered exactly as the article's lines are written, so no omission receipt of any kind accompanies this package. Citations: the retained excerpts carry no citation command at all, so no bibliography entry of the article is transcribed and no reference is redistributed. Reference adaptation: none was needed and none was made; every reference inside the delivered unit resolves inside it, so no unresolved reference or citation is produced. Printed slips kept literal: no printed word, symbol, hypothesis or number of the counted statement or of its proof was altered. Layout and class: the article's own class is \texttt{article} with packages including inputenc, babel, geometry, amsmath, amsfonts, amssymb, amsthm, bm, hyperref, kpfonts, enumerate, esint, mathalpha, todonotes; the wrapper is the lane's pinned \texttt{amsart} editorial wrapper at 10pt on A4, which restates the theorem-like environments the retained text needs and adds the article's own macro definitions that the retained text needs (\texttt{\textbackslash RR}, \texttt{\textbackslash divg}, \texttt{\textbackslash faskip}, \texttt{\textbackslash loc}, \texttt{\textbackslash supp}), with extra wrapper packages (bm, bbm, dsfont, xspace, cancel, mathtools). The delivered numbering is the unit's own. Assets: no retained span contains a figure environment, a graphics-inclusion command, a PDF-page inclusion, a table or a TikZ picture; no image file is copied into this package, the package declares no asset dependency and no image file is redistributed with it. Licence: version arXiv:2605.03553v1 of this article is distributed under the Creative Commons Attribution 4.0 International licence, as recorded in the arXiv licence metadata for this exact version and shown on the abstract page https://arxiv.org/abs/2605.03553v1. A full rights notice is delivered at licenses/arxiv-2605.03553-CC-BY-4.0.txt. Characters: every retained excerpt is pure ASCII, so the delivered engine, XeLaTeX with its default Latin Modern text font, reports no missing character.
\begin{lem}
	\label{lem:mp}
	Let $\Omega \subset \RR^{n}$ be a smooth bounded open set, $G: \Omega \rightarrow \RR^{n \times n}$ is a bounded, uniformly elliptic matrix, and $f\in C(\Omega)$. Let $u \in H^{1}(\Omega) \cap C(\Omega)$, $v \in H^{1}(\Omega)$ be nonnegative functions satisfying 
	\[
		- \divg(G\nabla u) = f\mathbb{1}_{\{u>0\}}, ~~ -\divg(G\nabla v)\geq f ~~ \text{in } \Omega,
	\]
	in the weak sense, and $v\geq u$ on $\partial \Omega$. Then, $v\geq u$ in $\Omega$. 
\end{lem}

\begin{equation}
\label{eq:obstacle_f}
	\begin{cases}
		-\Delta v = \left(\overline{\alpha} - f(y)  \right) \mathbb{1}_{\{v > 0\}}, \\
		v \geq 0, v\in L^{1}(\RR^{2}).
	\end{cases}
\end{equation}

\begin{prop}
	\label{prop:existence}
	Given $ \overline{\alpha} > 0 $, there exist a unique weak solution $ v \in H^{1}(\RR^{2}) $ to \eqref{eq:obstacle_f}.
\end{prop}

\begin{prop}
	\label{prop:nondeg_q_bound}
	$ \limsup_{n\rightarrow \infty} \frac{\beta_n}{\varepsilon_n^{\gamma}} < \infty$
\end{prop}

\begin{equation}
\label{eq:obst_noncoercive}
\begin{cases}
	-\Delta v = \left( \overline{\alpha} - x_1^{4} - x_1^{2} x_2^{2} \right) \mathbb{1}_{\{v>0\}} & \text{ in }\RR^{2},\\
	v\geq 0, v\in L^{1}(\RR^{2}). &
\end{cases}
\end{equation}

\begin{prop}
	\label{prop:exist_noncoercive}
	There exists a unique solution $ v \in H^{1}(\RR^{2}) $ to \eqref{eq:obst_noncoercive}.
\end{prop}

\begin{proof}
	The uniqueness part follows by the same argument as in the proof of Proposition \ref{prop:existence}.

	For $ \delta \in (0,1) $, define the problem
	\begin{equation}
	\label{eq:min_delta}
		\min \left\{ J_{\delta}(v) := \int_{\RR^{2}} \frac{|\nabla v|^{2}}{2} + \left( x_1^{4} + x_1^{2}x_2^{2} + \delta x_2^{4} \right)  v \mathrm{d}y: v\in H^{1}(\RR^{2}),~ v\geq 0, ~ \int_{\RR^{2}}v = 1  \right\}.
	\end{equation}

	We conclude that $ v_{\delta} $ is the unique minimizer of \eqref{eq:min_delta}. Hence, it satisfies
	\begin{equation}
		\label{eq:obstacle_delta}
		\begin{cases}
			-\Delta v_{\delta} = \left( \alpha_{\delta} - x_1^{4} - x_1^{2}x_2^{2} - \delta x_2^{4}	\right)\mathbb{1}_{\{v_{\delta} > 0 \}} & \text{ in }\RR^{2},\\
			v_{\delta} \geq 0, \int_{\RR^{2}}v_{\delta} = 1. &
		\end{cases}
	\end{equation}
	
	Suppose $ \limsup_{\delta \rightarrow 0}\alpha_{\delta} = +\infty $. Replicating the argument in the proof of Proposition \ref{prop:nondeg_q_bound}, we get 
	\begin{equation}
		\label{eq:inf_noncoercive}
		v_{\delta} \geq C\alpha_{\delta} ~\text{ in }B_r,
	\end{equation}
for some $ r > 0 $ independent of $ \delta $. This implies $ \|v_{\delta} \|_{L^{1}(\RR^{2})} > 1$ for some $ \delta $. This is a contradiction, therefore there is some $ C_0 $ such that $ \alpha_{\delta} \leq C_0 $ for all $ \delta \in (0,1) $. 

	Consider the function
	\[\varphi(x) = \begin{cases}
		\alpha_{\delta}^{3/2} \left( \dfrac{\sqrt{5}}{3} - \dfrac{1}{2} \left( \dfrac{x_1}{\alpha_{\delta}^{1/4}} \right)^{2} +\dfrac{1}{30} \left( \dfrac{x_1}{\alpha_{\delta}^{1/4}} \right)^{6}\right) & \text{ if }  \dfrac{x_1^{4}}{\alpha_{\delta}} \leq 5, \\
		0 & \text{ otherwise,}
	\end{cases}\]
which solves 
	\[
		\begin{cases}
			-\Delta \varphi = \left( \alpha_{\delta} - x_1^{4} \right) \mathbb{1}_{\{\varphi > 0\}} \faskip \text{in } \RR^{2},
			\varphi \geq 0 & \text{in }\RR^{2}.
		\end{cases}
	\]
	In particular, $ -\Delta \varphi \geq \alpha_{\delta} - x_1^{4} - x_1^{2}x_2^{2} - (\delta x_2)^{6} $ in $ \RR^{2} $. Since $ v_{\delta} \equiv 0 $ on $ \partial B_{\delta^{-1} R} $, we can use Lemma \ref{lem:mp} in the ball $ B_{\delta^{-1}R} $ to get $ v_{\delta} \leq \varphi $ in $ B_{\delta^{-1}R} $. 

	This yields that 
	\[
		v_{\delta} \leq \frac{\sqrt{5}\alpha_{\delta}^{3/2}}{3} \faskip \text{ in } \RR^{2},
	\]
	which, together with \eqref{eq:inf_noncoercive}, implies that $ \liminf_{\delta \rightarrow 0} \alpha_{\delta} > 0 $. Moreover,
	\[
		\supp (v_{\delta})\subset (-(5C_0)^{1/4}, (5C_0)^{1/4}) \times \RR \faskip \forall \delta \in (0,1).
	\]

	We can thus take a subsequence $ \delta_n \rightarrow 0 $ such that $ \alpha_{n} \rightarrow \alpha > 0$. Moreover, the right-hand side in \eqref{eq:obstacle_delta} is bounded in $ L^{\infty} $, so that, up to a further subsequence, $ v_{\delta_n} \rightharpoonup \overline{v} $ weakly in $ W^{2,p}_{\loc}(\RR^{2}) $, with $ \overline{v} \in L^{1}(\RR^{2}) $. This implies that 
	\[ 
		v(x) = \left( \frac{\overline{\alpha}}{\alpha} \right)^{2/3} \overline{v}\left( \left( \frac{\alpha}{\overline{\alpha}} \right)^{1/4} x \right)
	\]
solves \eqref{eq:obst_noncoercive}.
\end{proof}

\end{document}
